\begin{abstract}
Mathematical reasoning remains a formidable obstacle for language models, largely due to its inherent complexity and structured demands. In this work, we present DeepSeekMath~7B, an extension of the DeepSeek-Coder-Base-v1.5 7B model that undergoes continued pre-training with a corpus of 120B math-specific tokens collected from Common Crawl, integrated with both natural language and programming data. DeepSeekMath~7B secures an outstanding 51.7\% accuracy on the challenging MATH competition benchmark, attaining results comparable to those of Gemini-Ultra and GPT-4—all without external toolkits or ensemble voting. When leveraging self-consistency across 64 sampled outputs, DeepSeekMath~7B achieves 60.9\% on MATH. The exceptional mathematical reasoning shown by DeepSeekMath~is attributable to two main innovations: firstly, an advanced data selection pipeline that taps into the rich resources available in web-based datasets; and secondly, the development of Group Relative Policy Optimization (GRPO), an adaptation of Proximal Policy Optimization (PPO), which simultaneously strengthens reasoning performance and reduces PPO's memory demands.
\end{abstract}

\section{Introduction}

The emergence of large language models (LLMs) has transformed mathematical reasoning methodologies within artificial intelligence, leading to notable progress on quantitative benchmarks such as \citep{MATH} and challenges in geometric reasoning such as \citep{trinh2024solving}. Beyond these benchmarks, LLMs have become vital assets for human mathematicians addressing sophisticated mathematical questions \citep{tao}. Nonetheless, high-performing systems like GPT-4 \citep{gpt4} and Gemini-Ultra \citep{gemini} remain proprietary, while the state-of-the-art open-source alternatives continue to lag considerably in effectiveness.

This paper presents DeepSeekMath, a specialized language model tailored for mathematical reasoning, which substantially surpasses the performance of existing open-source systems and closely rivals GPT-4 on rigorous academic benchmarks. Central to this advance is the construction of the DeepSeekMath Corpus, a high-caliber pre-training resource containing 120B tokens specifically focused on mathematics. This corpus is derived from Common Crawl (CC) via a classifier built on fastText \citep{joulin2016fasttext}. The initial classifier iteration is trained with OpenWebMath \citep{openwebmath} data as positive samples and a varied assortment of non-mathematical web content as negative samples. Following this, the classifier is utilized to extract further positive candidates from CC, which are subsequently subjected to human validation for greater precision. An updated classifier is then trained using this enriched data, enhancing classification efficacy. Empirical results show the superior quality of the resulting dataset, as evidenced by the DeepSeekMath-Base 7B model, which attains 64.2\% on GSM8K \citep{gsm8k} and 36.2\% on the MATH competition dataset \citep{MATH}, outperforming models like Minerva 540B \citep{minerva}. Furthermore, owing to its multilingual design, the DeepSeekMath Corpus facilitates marked gains on Chinese mathematics benchmarks \citep{wei2023cmath,agieval}. We contend that our methodology in mathematical data preparation marks an initial step for broader research and anticipate considerable potential for further progress.

DeepSeekMath-Base is constructed by initializing with DeepSeek-Coder-Base-v1.5 7B \citep{deepseek-coder}, since empirical evidence suggests that models pre-trained on code exhibit better performance than those beginning from general-purpose language models. In addition, our findings indicate that mathematical training not only augments mathematical competencies but also bolsters overall reasoning performance, as demonstrated by improvements on MMLU \citep{mmlu} and BBH \citep{bbh} benchmarks.

Subsequent to pre-training, DeepSeekMath-Base undergoes mathematical instruction tuning utilizing datasets focused on chain-of-thought \citep{cot}, program-of-thought \citep{pot,pal}, and tool-integrated reasoning \citep{tora}. The resultant model, DeepSeekMath-Instruct 7B, outperforms all existing 7B-scale peers and achieves parity with open-source instruction-tuned models at the 70B parameter scale.

Moreover, we present Group Relative Policy Optimization (GRPO), an alternative reinforcement learning (RL) approach based on Proximal Policy Optimization (PPO) \citep{schulman2017proximal}. Unlike traditional PPO, GRPO dispenses with a critic model, employing group-based score estimates for the baseline, which leads to significant reductions in computational requirements. Leveraging only a targeted portion of English instruction tuning data, GRPO achieves marked gains over DeepSeekMath-Instruct in both in-domain (GSM8K: 82.9\% $\rightarrow$ 88.2\%, MATH: 46.8\% $\rightarrow$ 51.7\%) and out-of-domain (e.g., CMATH: 84.6\% $\rightarrow$ 88.8\%) mathematical evaluations during RL training. Additionally, we develop a unified framework for understanding a variety of methods—including Rejection Sampling Fine-Tuning (RFT) \citep{yuan2023scaling}, Direct Preference Optimization (DPO) \citep{dpo}, PPO, and GRPO—by viewing them as either direct or streamlined RL strategies. Comprehensive experiments are conducted, covering comparisons such as online versus offline learning, outcome versus process supervision, and single-turn versus iterative RL, with the aim of dissecting core components of the unified framework. Finally, we analyze the mechanisms by which our RL strategy enhances instruction-tuned model performance and outline future avenues for advancing RL effectiveness within this framework.

\subsection{Contributions}
Our work advances both scalable math pre-training and the study and evaluation of reinforcement learning techniques.

\textbf{Math Pre-Training at Scale}

\begin{itemize}[topsep=0pt]
    \item We demonstrate that Common Crawl, an openly available web corpus, contains significant mathematical resources. By constructing a rigorous data selection framework, we introduce the DeepSeekMath Corpus—a curated dataset containing 120B tokens sourced from web content with mathematical relevance. This dataset is nearly 7 times the scale of the math-specific web pages in Minerva \citep{minerva}, and approximately 9 times larger than OpenWebMath \citep{openwebmath}.

    \item Our DeepSeekMath-Base 7B model, following pre-training, delivers performance on par with the substantially larger Minerva 540B \citep{minerva}. This suggests that sheer parameter count does not exclusively determine a model's mathematical reasoning proficiency. It further indicates that smaller models, when trained on carefully filtered data, can excel at mathematical tasks.

    \item Through extensive math pre-training trials, we observe that conducting code pre-training prior to math-specific training enhances a model’s ability to tackle mathematical problems, regardless of the availability of computational tools. This provides partial support for the notion that code pre-training augments reasoning capacity, specifically in mathematical domains.

    \item Despite the prevalent use of arXiv documents for training in math-focused research, incorporating arXiv content leads to no discernible gains across all math benchmarks considered in this study.
\end{itemize}

\textbf{Exploration and Analysis of Reinforcement Learning}

\begin{itemize}[topsep=0pt]
    \item We present Group Relative Policy Optimization (GRPO), a reinforcement learning framework that offers both efficacy and efficiency. By bypassing the need for a critic and instead leveraging group-based baseline estimation, GRPO dramatically reduces the computational demands associated with traditional algorithms like Proximal Policy Optimization (PPO).
    \item Our experiments establish that GRPO substantially boosts the capabilities of our instruction-tuned system, DeepSeekMath-Instruct, relying exclusively on instruction-tuning datasets. Additionally, the application of GRPO results in observable improvements in generalization to domains beyond the training distribution.
    \item We introduce a cohesive framework for interpreting and relating various approaches, including RFT, DPO, PPO, and GRPO. In-depth experimental comparisons—covering online versus offline paradigms, outcome-focused versus process-oriented supervision, and single-step versus multi-step reinforcement learning—enable a thorough examination of this unified perspective.
    \item Leveraging our integrative paradigm, we analyze the mechanisms underlying the efficacy of reinforcement learning for language models and outline several promising research directions aimed at enhancing reinforcement learning outcomes in large language models.
\end{itemize}

\subsection{Summary of Evaluations and Metrics}

\begin{itemize}[topsep=0pt]
    \item \textbf{English and Chinese Mathematical Reasoning}: We systematically evaluate our models across both English and Chinese mathematical reasoning tasks, spanning a spectrum from elementary through collegiate problem sets. English language assessments utilize benchmarks such as GSM8K \citep{gsm8k}, MATH \citep{MATH}, SAT \citep{llemma}, OCW Courses \citep{minerva}, and MMLU-STEM \citep{mmlu}, while Chinese benchmarks include MGSM-zh \citep{mgsm}, CMATH \citep{wei2023cmath}, Gaokao-MathCloze \citep{agieval}, and Gaokao-MathQA \citep{agieval}. We measure the proficiency of the models in formulating independent written solutions (without the aid of external tools) as well as their effectiveness in problem-solving with Python.

    On English mathematical benchmarks, DeepSeekMath-Base matches the performance of the proprietary Minerva 540B \citep{minerva} and consistently outperforms open-source base models—such as Mistral 7B \citep{mistral} and Llemma-34B \citep{llemma}—irrespective of previous exposure to mathematics-specific pre-training, often by considerable margins. Distinctly, DeepSeekMath-Base exhibits even stronger results on Chinese evaluations, a trend attributable to our decision to go beyond the English-only datasets utilized in prior studies \citep{minerva, llemma}, integrating high-quality data from a broader linguistic scope. Enhanced by instruction tuning and reinforcement learning, DeepSeekMath-Instruct and DeepSeekMath-RL achieve robust outcomes, notably surpassing the 50\% accuracy threshold on the highly challenging, competition-grade MATH dataset—a milestone within the open-source community.
\end{itemize}

\item \textbf{Formal Mathematics}: DeepSeekMath-Base is assessed on the informal-to-formal theorem proving benchmark from \citep{dsp_proof} using the miniF2F dataset \citep{minif2f} and Isabelle \citep{isabelle} as the designated proof assistant. The results indicate that DeepSeekMath-Base achieves robust performance in few-shot autoformalization tasks.
\item \textbf{Natural Language Understanding, Reasoning, and Code}: To comprehensively evaluate the model’s proficiency in general understanding, reasoning, and coding, we test DeepSeekMath-Base on several benchmarks: Massive Multitask Language Understanding (MMLU) \citep{mmlu} featuring 57 multiple-choice subject areas, BIG-Bench Hard (BBH) \citep{bbh} which includes 23 highly challenging, multi-step reasoning tasks, along with the standard coding benchmarks HumanEval \citep{codex} and MBPP \citep{mbpp}. Pre-training on mathematical data is shown to improve both language comprehension and reasoning abilities.
\end{itemize}

\section{Math Pre-Training}

\subsection{Data Collection and Decontamination} 
Here, we describe the methodology for constructing the DeepSeekMath Corpus using data sourced from Common Crawl. Figure \ref{fig:data_collect} illustrates the iterative process utilized to systematically aggregate a large-scale mathematical dataset from Common Crawl, commencing from a small, high-quality seed corpus comprised of math-related data. Notably, this framework is extendable to other specialized domains, such as code.

For the initial seed, we utilize OpenWebMath \citep{openwebmath}, recognized for its quality mathematical web texts. With this as the base, we train a fastText model \citep{joulin2016fasttext} to retrieve further web pages exhibiting characteristics similar to OpenWebMath. In particular, we randomly extract 500,000 instances from the seed corpus as positive samples, and select an equal number (500,000) of web pages from Common Crawl to serve as negative samples. Training is conducted using a publicly available fastText library\footnote{\url{https://fasttext.cc}}, configured with a vector dimensionality of 256, a learning rate of 0.1, a maximum n-gram length of 3, a minimum occurrence threshold of 3 for each word, and three epochs. To minimize the redundancy in the original Common Crawl dataset, we apply URL deduplication as well as near-deduplication strategies, which reduces the collection to 40 billion HTML web pages. Using the trained fastText model, we identify mathematical pages from the deduplicated corpus. We then prioritize these pages according to the relevance scores computed by fastText, retaining only those with the highest scores to eliminate lower-quality material. The retained data volume is further evaluated through pre-training experiments across varying scales: 40B, 80B, 120B, and 160B tokens. In this initial cycle, we opt to preserve the top 40B tokens.

Following the initial round of data collection, a significant number of mathematical web pages remain uncollected. This issue arises chiefly because the fastText model’s set of positive training examples lacks adequate variety. To address this, we broaden the seed corpus by integrating additional mathematical web sources, thereby facilitating improved training of the fastText model. Our approach begins by partitioning the entirety of Common Crawl into mutually exclusive domains, where each domain is comprised of web pages sharing the same base URL. For every domain, we compute the proportion of pages collected during the initial round. Any domain in which more than 10\% of pages have been gathered is designated as math-related (for instance, \url{mathoverflow.net}).

After identifying such domains, we conduct manual annotation of URLs that are likely to contain mathematical material within these domains (such as \url{mathoverflow.net/questions}). Uncollected web pages under these annotated URLs are incorporated into our seed set. This methodology enables us to aggregate additional positive samples, thus enhancing the fastText model’s capacity to identify mathematical content during subsequent rounds of data collection. Over four iterative rounds, our process yields a final dataset comprising 35.5M mathematical web pages, encompassing 120B tokens in total. Notably, by the fourth round, we observe that approximately 98\% of all data had already been acquired during the third iteration, prompting us to terminate further collection.

In order to prevent contamination of evaluation benchmarks, we apply the filtering protocol outlined in \cite{deepseek-coder}, systematically excluding web pages containing any questions or answers from both English mathematical benchmarks (e.g., GSM8K \citep{gsm8k}, MATH \citep{MATH}) and Chinese mathematical benchmarks (e.g., CMATH \citep{wei2023cmath}, AGIEval \citep{agieval}). The filtering mechanism involves removing any segment that has an exact 10-gram match with any substring present in the benchmark corpora. For benchmark samples shorter than 10-grams but with a length of at least 3-grams, precise matching is applied to exclude potentially contaminated web pages from the mathematical training corpus.

\subsection{Evaluating the DeepSeekMath~Corpus Quality}

To assess the efficacy of the DeepSeekMath~Corpus, we carry out pre-training experiments in comparison with several recently introduced mathematics-focused corpora:

\begin{itemize}[topsep=0pt]
    \item \textbf{MathPile} \citep{mathpile}: a composite multi-source collection (8.9B tokens), composed of material from textbooks, Wikipedia, ProofWiki, CommonCrawl, StackExchange, and arXiv, where over 85\% is drawn from arXiv;
    \item \textbf{OpenWebMath} \citep{openwebmath}: a CommonCrawl-derived dataset curated for mathematical relevance, consisting of 13.6B tokens;
    \item \textbf{Proof-Pile-2} \citep{llemma}: a mathematical dataset aggregating OpenWebMath, AlgebraicStack (10.3B tokens of mathematical code), and papers from arXiv (28.0B tokens). For experiments with Proof-Pile-2, we adopt the arXiv:Web:Code sampling ratio of 2:4:1, as recommended by \cite{llemma}.
\end{itemize}

\subsubsection{Training Setting}
\label{sec:quality-policy}
We fine-tune a 1.3B-parameter general-purpose language model, which is architecturally aligned with the DeepSeek LLMs \citep{deepseek-llm} and here referred to as DeepSeek-LLM 1.3B, on various mathematical datasets. Each model undergoes training exclusively on its respective mathematical corpus for 150B tokens. All experiments utilize the resource-efficient HAI-LLM \citep{haillm} training infrastructure. Adhering to DeepSeek LLMs’ established protocols, optimization is performed with AdamW \citep{adamW}, using parameters $\beta_1=0.9$, $\beta_2=0.95$, and $\mathrm{weight\_decay}=0.1$. The learning rate follows a staged schedule: it ramps up to its maximum after 2,000 warmup steps, drops to 31.6\% of the peak by 80\% of training, and further tapers to 10.0\% by the 90\% mark. The peak learning rate is capped at 5.3e-4, and models are trained using a batch size of 4 million tokens with sequences up to 4,000 tokens in length.

\subsubsection{Evaluation Results}

\textbf{DeepSeekMath~Corpus exhibits outstanding quality, comprehensive multilingual scope, and unparalleled scale.}

\begin{itemize}[topsep=0pt]
    \item \textbf{High-quality}:
    Downstream performance is assessed across 8 mathematical evaluation sets, utilizing few-shot chain-of-thought prompting \cite{cot}. According to the outcomes in Table \ref{tab:corpora_comparison}, models trained with DeepSeekMath~Corpus consistently outperform those trained on alternative datasets. Figure \ref{fig:corpora_comparisons} highlights that DeepSeekMath-trained models surpass those trained on Proof-Pile-2, even after just 50B tokens (one full pass through Proof-Pile-2), demonstrating that DeepSeekMath~Corpus delivers superior average sample quality.
    \item \textbf{Multilingual}:
    DeepSeekMath~Corpus contains mathematical material in several languages, with English and Chinese making up the bulk of the content. Table \ref{tab:corpora_comparison} shows that leveraging DeepSeekMath~Corpus significantly boosts mathematical reasoning abilities in both English and Chinese. In contrast, most previous corpora, which predominantly emphasize English, fail to substantially benefit non-English mathematical reasoning and, in some cases, negatively affect Chinese language performance.
    \item \textbf{Large-scale}:
    DeepSeekMath~Corpus surpasses prior mathematical corpora by a considerable margin in size. As visualized in Figure \ref{fig:corpora_comparisons}, the DeepSeek-LLM 1.3B trained on DeepSeekMath~Corpus displays a notably sharper learning trajectory and more persistent gains. In comparison, models trained on smaller, previously-used corpora quickly exhaust available unique data, leading to rapid convergence and limited further improvement.
\end{itemize}

\subsection{Training and Evaluating DeepSeekMath-Base 7B}

This section presents DeepSeekMath-Base 7B, a foundational model designed for robust reasoning capabilities, particularly in mathematical domains. The model begins with weights from DeepSeek-Coder-Base-v1.5 7B \citep{deepseek-coder} and is subsequently trained on a dataset comprising 500B tokens. The training data are distributed as follows: 56\% originates from the DeepSeekMath Corpus, 4\% from AlgebraicStack, 10\% from arXiv, 20\% from Github code repositories, and the remaining 10\% consists of natural language samples from Common Crawl, encompassing both English and Chinese text. The training approach largely mirrors the procedure described in Section \ref{sec:quality-policy}, except that the peak learning rate is adjusted to 4.2e-4, with the batch size set at 10M tokens.

A thorough evaluation of DeepSeekMath-Base 7B is conducted to gauge its mathematical problem-solving proficiency. This includes its capability to generate complete solutions independently, utilize external tools for mathematical problem solving, and engage in rigorous formal theorem proving. Furthermore, to better characterize the overall abilities of the base model, we report its performance on general tasks such as natural language comprehension, logical reasoning, and programming.

\paragraph{Mathematical Problem Solving with Step-by-Step Reasoning}

The problem-solving ability of DeepSeekMath-Base is assessed through few-shot chain-of-thought prompting \citep{cot} across eight distinct benchmarks presented in both English and Chinese. These datasets span quantitative reasoning tasks—including benchmarks like GSM8K \citep{gsm8k}, MATH \citep{MATH}, and CMATH \citep{wei2023cmath}—as well as multiple-choice assessments such as MMLU-STEM \citep{mmlu} and Gaokao-MathQA \citep{agieval}. The scope of these benchmarks covers a wide array of mathematical disciplines, ranging from foundational to advanced (elementary through college level).

As presented in Table \ref{tab:base_math_cot}, DeepSeekMath-Base 7B achieves superior results across all eight benchmarks when compared to other open-source base models, which include both the popular Mistral 7B \citep{mistral} and the newly introduced Llemma 34B \citep{llemma}, trained with mathematical data from Proof-Pile-2 \citep{llemma}. In particular, DeepSeekMath-Base establishes a new milestone on the MATH dataset, outperforming the best open-source base alternatives by more than 10 percentage points in absolute terms. Furthermore, it surpasses Minerva 540B \citep{minerva}—a proprietary model with 77 times the parameter count, built upon PaLM \citep{palm} and further refined with mathematical corpora.

\paragraph{Mathematical Problem Solving with Tool Use} For program-assisted mathematical reasoning, we assess models on GSM8K and MATH datasets by applying few-shot program-of-thought prompts \citep{pot,pal}. Each problem is solved by instructing the model to write a Python script, making use of external libraries such as \textit{math} and \textit{sympy} to manage complex calculations. The program's output is regarded as the solution. According to Table \ref{tab:base_math_pot_proof}, DeepSeekMath-Base 7B consistently outperforms the previous best model, Llemma 34B.

\paragraph{Formal Mathematics}

Formal proof automation is increasingly recognized for its potential to enhance mathematical proof reliability and efficiency. We benchmark DeepSeekMath-Base 7B on informal-to-formal proof generation as described in \citep{dsp_proof}. Here, the model receives an informal conjecture, its formal expression, and an informal proof, then attempts to synthesize a formal proof. Using the miniF2F suite \citep{minif2f}, which targets formalization of Olympiad-level problems, formal proofs in Isabelle are generated for each task with few-shot examples. In line with the approach in \cite{dsp_proof}, the model produces proof outlines which are completed by the Sledgehammer automated prover \citep{sledgehammer}. Table \ref{tab:base_math_pot_proof} reports that DeepSeekMath-Base 7B shows robust performance in automated proof formalization.

\paragraph{Natural Language Understanding, Reasoning, and Code}

For evaluation in natural language understanding, we consider MMLU \citep{mmlu}; for reasoning, BBH \citep{bbh}; and for programming, HumanEval \citep{codex} and MBPP \citep{mbpp}. Table \ref{tab:understanding_reasoning_code} shows that DeepSeekMath-Base 7B achieves substantial improvement on both MMLU and BBH over its predecessor, DeepSeek-Coder-Base-v1.5 \citep{deepseek-coder}, demonstrating that mathematical training enhances general language comprehension and reasoning. Furthermore, the integration of code tokens during ongoing training enables DeepSeekMath-Base 7B to sustain the coding capabilities of DeepSeek-Coder-Base-v1.5 across HumanEval and MBPP. Overall, DeepSeekMath-Base 7B achieves notably higher results than the general-purpose Mistral 7B \citep{mistral} across all three tasks involving reasoning and coding.

\section{Supervised Fine-Tuning}

\subsection{SFT Data Curation}

We developed a mathematical instruction-tuning dataset that incorporates both English and Chinese problems, encompassing a range of mathematical disciplines and difficulties. Each problem is provided with a solution in chain-of-thought (CoT) \citep{cot}, program-of-thought (PoT) \citep{pot,pal}, or tool-integrated reasoning format \citep{tora}, resulting in a total of 776K training instances.

\begin{itemize}[topsep=0pt]
    \item \textbf{English mathematical datasets}: We enhance GSM8K and MATH problems by supplementing them with tool-integrated solutions. Additionally, we include a selected portion of MathInstruct \citep{MathInstruct} and the training set from Lila-OOD \citep{lila}, where solutions employ either CoT or PoT strategies. Our English dataset spans a broad range of mathematical domains such as algebra, probability, number theory, calculus, and geometry.
    \item \textbf{Chinese mathematical datasets}: We gather K-12 mathematics questions in Chinese across 76 distinct sub-categories, like linear equations, and annotate the solutions using both CoT and tool-integrated reasoning frameworks.
\end{itemize}

\subsection{Training and Evaluating DeepSeekMath-Instruct 7B}

Here, we detail the process of instruction tuning DeepSeekMath-Instruct 7B, using DeepSeekMath-Base as its foundation. Training samples are concatenated at random, up to a maximum of 4K tokens per context window. The model undergoes 500 steps of training, using a batch size of 256, with the learning rate fixed at 5e-5.

The mathematical capabilities of the models are assessed, both in standard and tool-assisted contexts, using four quantitative reasoning benchmarks in English and Chinese. We conduct comparative evaluation with top-performing models available at the time:

\begin{itemize}[topsep=0pt]
    \item \textbf{Closed-source models} assessed include: (1) the GPT family, notably GPT-4 \citep{gpt4} and GPT-4 Code Interpreter~\footnote{\url{https://openai.com/blog/chatgpt-plugins\#\#code-interpreter}}, (2) Gemini Ultra and Pro \citep{gemini}, (3) Inflection-2 \citep{inflection-2}, (4) Grok-1~\footnote{\url{https://x.ai/model-card}}, as well as recent releases from Chinese developers including (5) Baichuan-3~\footnote{\url{https://www.baichuan-ai.com}}, (6) and the newest GLM-4~\footnote{\url{https://open.bigmodel.cn/dev/api\#glm-4}} within the GLM suite \citep{glm}.
\end{itemize}

These models are designed for general applications and typically have undergone multiple alignment processes. 
\item \textbf{Open-source models} consist of both general-purpose and mathematics-enhanced systems, including: (1) DeepSeek-LLM-Chat 67B \citep{deepseek-llm}, (2) Qwen 72B \citep{qwen}, (3) SeaLLM-v2 7B \citep{seallm}, and (4) ChatGLM3 6B \citep{chatglm3}, alongside models optimized for mathematical capabilities such as (5) InternLM2-Math 20B~\footnote{\url{https://github.com/InternLM/InternLM-Math}}, which extends InternLM2 via math-oriented pretraining and subsequent instruction fine-tuning; (6) Math-Shepherd-Mistral 7B, incorporating PPO training \citep{schulman2017proximal} on top of Mistral 7B \citep{mistral} with a process-based reward model; (7) WizardMath series \citep{wizardmath}, which enhances math reasoning in Mistral 7B and Llama-2 70B \citep{llama2} using evolve-instruct (an instruction tuning technique with AI-generated prompts) and PPO training, leveraging data predominantly from GSM8K and MATH; (8) MetaMath 70B \citep{metamath}, which fine-tunes Llama-2 70B on an extended version of GSM8K and MATH; (9) ToRA 34B \cite{tora}, built by further training CodeLlama 34B for mathematical reasoning with integrated tool use; and (10) MAmmoTH 70B \citep{MathInstruct}, an instruction-tuned variant of Llama-2 70B on MathInstruct.

Table \ref{tab:sft_rl_math} presents results under a strict evaluation regime prohibiting external tool usage, where DeepSeekMath-Instruct 7B excels in systematic reasoning steps. Remarkably, this model achieves at least 9\% higher accuracy than every other open-source and most proprietary models (such as Inflection-2 and Gemini Pro) on the advanced MATH dataset. This superiority holds even compared to considerably larger models (like Qwen 72B) and those benefiting from math-oriented RL techniques (e.g., WizardMath-v1.1 7B). Although DeepSeekMath-Instruct is on par with proprietary Chinese models like GLM-4 and Baichuan-3 on MATH, it still falls short relative to GPT-4 and Gemini Ultra.

When assessed under a paradigm allowing both natural language reasoning and the deployment of programmatic tools, DeepSeekMath-Instruct 7B nearly achieves a 60\% accuracy rate on MATH, outperforming all other open-source systems to date. On additional benchmarks, its results are competitive with DeepSeek-LLM-Chat 67B, the previous best-performing model, which is tenfold larger in scale.

\section{Reinforcement Learning}

\subsection{Group Relative Policy Optimization} Reinforcement learning (RL) has demonstrated its efficacy in further strengthening the mathematical reasoning capabilities of large language models (LLMs) post-Supervised Fine-Tuning (SFT) \citep{wang2023math,wizardmath}. In what follows, we present our RL framework, Group Relative Policy Optimization (GRPO), which is both efficient and effective.

\subsubsection{From PPO to GRPO} Proximal Policy Optimization (PPO) \citep{schulman2017proximal} is a mainstream actor-critic algorithm commonly employed to refine LLMs in the RL fine-tuning phase \citep{ouyang2022training}. The method aims to optimize LLM policies through the following surrogate objective:

\begin{equation}
\footnotesize
    \mathcal{J}_{PPO}(\theta) = \mathbb{E}{[q \sim P(Q), o \sim \pi_{\theta_{old}}(O|q)]} \frac{1}{|o|} \sum_{t=1}^{|o|} \min \left[ \frac{\pi_\theta(o_{t} | q, o_{<t})}{\pi_{\theta_{old}}(o_{t} | q, o_{<t})} A_{t}, \text{clip} \left( \frac{\pi_\theta(o_{t} | q, o_{<t})}{\pi_{\theta_{old}}(o_{t} | q, o_{<t})}, 1 - \varepsilon, 1 + \varepsilon \right)  A_{t

In this context, $\pi_{\theta}$ and $\pi_{\theta_{old}}$ denote the present and preceding policy models, while $q$ and $o$ refer to questions and their respective outputs sampled either from the question corpus or from the old policy $\pi_{\theta_{old}}$. The parameter $\varepsilon$ is a clipping hyper-parameter utilized by PPO to promote more stable training dynamics. The advantage term $A_t$ is calculated via Generalized Advantage Estimation (GAE) \citep{gae}, leveraging both the reward sequence $\{r_{\ge t}\}$ and a value function $V_{\psi}$, which is learned. Accordingly, PPO requires simultaneous training of a value function with the policy model. To counteract excessive exploitation of the reward model, a common technique incorporates a per-token KL divergence penalty into the reward at every token, referencing an external model as described by \citet{ouyang2022training}:

\begin{equation}
    r_{t} = r_\varphi(q, o_{\le t}) - \beta \log\frac{\pi_{\theta}(o_{t}|q, o_{<t})}{\pi_{ref}(o_{t}|q, o_{<t})},
\label{eq:PPO-reward}
\end{equation}

Here, $r_\varphi$ designates the reward model, $\pi_{ref}$ is the reference policy—typically the SFT initialization—and $\beta$ weights the KL penalty.

Given that the value function in PPO is generally another network with similar scale to the main policy model, this setup imposes heavy computational and memory requirements. During reinforcement learning, the value function serves as a baseline for the advantage estimate, primarily to reduce variance. However, in the case of large language models (LLMs), rewards are usually only assigned to the final token by the reward model, making accurate value prediction at every token challenging. To alleviate this difficulty, we introduce Group Relative Policy Optimization (GRPO) as depicted in Figure \ref{fig:grpo}. Unlike PPO, GRPO removes the necessity for value function approximation by substituting it with the mean reward across several outputs generated for the same input question, which serves as the baseline. Concretely, for each $q$, GRPO draws a collection of outputs $\{o_1, o_2, \cdots, o_G\}$ sampled from $\pi_{\theta_{old}}$, and the policy is updated by optimizing the objective function:

\begin{equation}
\footnotesize
\begin{split}
    \mathcal{J}_{GRPO}(\theta) &= \mathbb{E}{[q \sim P(Q), \{o_i\}_{i=1}^G \sim \pi_{\theta_{old}}(O|q)]}  \\
    & \frac{1}{G}\sum_{i=1}^G\frac{1}{|o_i|} \sum_{t=1}^{|o_i|} \left\{ \min \left[ \frac{\pi_\theta(o_{i,t} | q, o_{i,<t})}{\pi_{\theta_{old}}(o_{i,t} | q, o_{i,<t})} \hat{A}_{i,t}, \text{clip} \left( \frac{\pi_\theta(o_{i,t} | q, o_{i,<t})}{\pi_{\theta_{old}}(o_{i,t} | q, o_{i,<t})}, 1 - \varepsilon, 1 + \varepsilon \right)  \hat{A}_{i,t} \right] - \beta \mathbb{D}_{KL}\left[\pi_{\theta} || \pi_{ref}\right]\right\} ,
\end{split}
\label{eq:GRPO-obj}
\end{equation}

In this formulation, the hyper-parameters $\varepsilon$ and $\beta$ control the respective clipping and regularization strengths, and $\hat{A}_{i,t}$ is the advantage, which is now determined solely based on the relative rewards among the grouped outputs for a particular question; further explanation follows in the subsequent subsections. This group-relative calculation of advantage harmonizes with the reward model’s design, which is usually optimized using datasets containing pairwise or comparative evaluations over outputs addressing identical queries. Moreover, rather than integrating the KL regularization directly into the reward as in PPO, GRPO penalizes KL divergence by introducing it explicitly as a loss term, simplifying the computation of $\hat{A}_{i,t}$. Notably, in contrast to the KL penalty in Eq. (\ref{eq:PPO-reward}), the KL divergence in GRPO is estimated with an unbiased estimator from \citet

\begin{algorithm}[t]
  \small
  \caption{Iterative Group Relative Policy Optimization}
  \textbf{Input} initial policy model $\pi_{\theta_{\text{init}}}$; reward models $r_\varphi$; task prompts $\mathcal{D}$; 
  hyperparameters $\varepsilon$, $\beta$, $\mu$
  \begin{algorithmic}[1]
    \State Set the policy model $\pi_\theta$ to $\pi_{\theta_{\text{init}}}$
    \For{iteration = 1, \dots, I}
       \State Set the reference model $\pi_{ref}$ equal to $\pi_{\theta}$
      \For{step = 1, \dots, M}
      \State Draw a minibatch $\mathcal{D}_b$ from $\mathcal{D}$
      \State Copy current policy to $\pi_{\theta_{old}} \leftarrow \pi_{\theta}$ 
      \State For each question $q \in \mathcal{D}_b$, generate $G$ responses $\{o_i\}_{i=1}^G \sim \pi_{\theta_{old}} (\cdot \mid q) $
      \State Evaluate each output $o_i$ using $r_{\varphi}$ to produce rewards $\{r_i\}_{i=1}^{G}$
      \State For each $t$-th token of $o_i$, calculate $\hat{A}_{i,t}$ via group relative advantage estimation
      \For{GRPO iteration = 1, \dots, $\mu$}
        \State Adjust the policy model $\pi_{\theta}$ to maximize the GRPO objective (Equation \ref{eq:GC-GRPO})
      \EndFor
    \EndFor \State Update $r_\varphi$ incrementally through a replay-based approach.
    \EndFor 
  \end{algorithmic}
  \textbf{Output} $\pi_\theta$
  \label{alg:iter-grpo}
\end{algorithm}

\subsubsection{Outcome Supervision RL with GRPO} In this framework, given a question $q$, the previous policy model $\pi_{\theta_{old}}$ is used to sample a collection of responses $\{o_1, o_2, \cdots, o_G\}$. Each of these outputs receives an individual reward from a reward model, resulting in $G$ scores $\mathbf{r}=\{r_1, r_2, \cdots, r_G\}$. These rewards are standardized by removing the mean of the group and dividing by the standard deviation. Under outcome supervision, this normalized reward is attributed to the last token of each output $o_i$ and also uniformly assigned as the advantage $\hat{A}_{i, t}$ for all tokens, such that $\hat{A}_{i, t} = \widetilde{r}_i = \frac{r_i- {\rm mean}(\mathbf{r})}{{\rm std}(\mathbf{r})}$. The learning process then updates the policy to maximize the target in equation (\ref{eq:GRPO-obj}).

\subsubsection{Process Supervision RL with GRPO} Since outcome supervision limits feedback to the end of each response, it can be inadequate for supervising agents on intricate mathematical reasoning. Building on the methodology in \cite{wang2023math}, process supervision extends the approach by granting rewards at the conclusion of each reasoning step. Concretely, for a given question $q$ and generated batch $\{o_1, o_2, \cdots, o_G\}$, a process-based reward model evaluates each intermediate step, producing a nested set of rewards: $\mathbf{R} = \{ \{r_1^{index(1)},\cdots,r_1^{index(K_1)}\}, \cdots,  \{r_G^{index(1)},\cdots,r_G^{index(K_G)}\} \}$, with $index(j)$ denoting the position of the terminal token for step $j$ and $K_i$ marking the step count of the $i$-th output. The normalization of rewards follows by standardizing with the mean and standard deviation: $\widetilde{r}_i^{index(j)} = \frac{r_i^{index(j)} - {\rm mean(\mathbf{R})}}{{\rm std(\mathbf{R})}}$. The process supervision mechanism then sets the advantage for the $t$-th token as the sum over all normalized subsequent step rewards, i.e., $\hat{A}_{i, t} = \sum_{index(j) \ge t} \widetilde{r}_i^{index(j)}$, and

\subsubsection{Iterative RL with GRPO} Over the course of reinforcement learning, previously trained reward models may become inadequate for effectively guiding the updated policy model. To address this, we investigate iterative RL employing GRPO. As detailed in Algorithm \ref{alg:iter-grpo}, this iterative procedure involves generating new datasets for the reward model using samples produced by the evolving policy model, while continually updating the reward model itself with a replay buffer that integrates 10\% of past data. The reference model is set to the latest policy model, which is then further optimized using the recently updated reward model.

\subsection{Training and Evaluating DeepSeekMath-RL}

Our RL experiments are based on DeepSeekMath-Instruct 7B. The RL phase utilizes training examples in the chain-of-thought format drawn from GSM8K and MATH datasets within the SFT data, comprising roughly 144K examples. We intentionally omit SFT questions from other sources to assess how RL affects performance on benchmarks absent from the RL training stage. Reward model training data is prepared as outlined in \citep{wang2023math}. The initial reward model is trained from DeepSeekMath-Base 7B using a learning rate of 2e-5. For GRPO, the policy model adopts a learning rate of 1e-6, with the KL penalty coefficient fixed at 0.04. For each problem, we generate $64$ candidate completions. Sequence length is capped at 1024 tokens, and the batch size is set to 1024. After each exploration cycle, the policy model undergoes a single update step. Evaluation of DeepSeekMath-RL 7B mirrors the methodology used for DeepSeekMath-Instruct 7B. For DeepSeekMath-RL 7B, chain-of-thought-style tasks from GSM8K and MATH represent in-domain evaluation, while the remaining benchmarks serve as out-of-domain tasks.

Table \ref{tab:sft_rl_math} presents a comparative analysis of both open- and closed-source models employing chain-of-thought and tool-augmented reasoning across English and Chinese datasets. The results indicate: 1) DeepSeekMath-RL 7B achieves 88.2\% and 51.7\% accuracy on GSM8K and MATH benchmarks, respectively, when utilizing chain-of-thought reasoning. These scores outperform all open-source counterparts within the 7B–70B parameter range and exceed the capabilities of most closed-source systems. 2) Notably, DeepSeekMath-RL 7B’s training is restricted to chain-of-thought-format instruction tuning data from GSM8K and MATH, and it builds upon DeepSeekMath-Instruct 7B. Despite this narrow training focus, it consistently outperforms DeepSeekMath-Instruct 7B on all metrics, underscoring the advantages of reinforcement learning.

\section{Discussion}
This section details insights gained from both pre-training and RL experiments.

\subsection{Lessons Learnt in Pre-Training}

We begin by outlining our observations from the pre-training process. Unless otherwise stated, the experimental conditions mirror those described in Section \ref{sec:quality-policy}. For the purpose of this discussion, the term DeepSeekMath Corpus refers to the 89B-token dataset compiled during the second round of data collection.

\subsubsection{Code Training Benefits Mathematical Reasoning}

An oft-cited but rarely substantiated belief is that exposure to code during training enhances a model’s reasoning skills. Here, we offer empirical evidence supporting this claim, especially within mathematics: code-centric pre-training contributes positively to a model’s mathematical reasoning capabilities, both with and without auxiliary tools.

To systematically evaluate the effect of code-based training on mathematical reasoning, we conducted comparative studies using both a two-stage training regime and a single-stage training protocol:

\textbf{Two-Stage Training}

\begin{itemize}[topsep=0pt]
    \item \textbf{Code Training for 400B Tokens $\rightarrow$ Math Training for 150B Tokens}: DeepSeek-LLM 1.3B is initially trained with 400B code tokens, after which it undergoes further training on 150B tokens consisting of math data;
    \item \textbf{General Training for 400B Tokens $\rightarrow$ Math Training for 150B Tokens}: As a comparison, we also conduct training with 400B general tokens, sampled from a large-scale corpus compiled by DeepSeek-AI, before the same phase of math-focused training. This control is intended to evaluate whether code tokens offer specific benefits for mathematical reasoning over general tokens.
\end{itemize}

\textbf{One-Stage Training}

\begin{itemize}[topsep=0pt]
    \item \textbf{Math Training for 150B Tokens}: DeepSeek-LLM 1.3B is trained directly on 150B tokens of mathematical content without any pretraining on code or general data;
    \item \textbf{Training on a mixture of 400B Code Tokens and 150B Math Tokens}: Since the sequential math training following code training is observed to decrease coding proficiency, we also explore a setup where code and math tokens are intermixed for the entirety of the training in a single stage. This approach allows us to study if such mixing sustains improvements in mathematical reasoning and counters catastrophic forgetting.
\end{itemize}

\paragraph{Results}
Table \ref{tab:code-math} and Table \ref{tab:code-math-general-eval} report performance across the evaluated training paradigms.

Engagement with code data positively impacts program-aided mathematical reasoning, both in multi-stage and single-stage regimes. As detailed in Table \ref{tab:code-math}, the incorporation of code during the first stage alone markedly improves proficiency in solving tasks from GSM8K and MATH via Python-based methods. Incorporating an additional math-specific training stage leads to even greater gains. Notably, when code and math tokens are mixed throughout one-stage training, the catastrophic forgetting commonly seen in multi-stage training is considerably reduced, while coding abilities (Table \ref{tab:code-math-general-eval}) and program-aided mathematical performance (Table \ref{tab:code-math}) are also enhanced.

Furthermore, training on code benefits mathematical reasoning independent of tool usage. Two-stage training, beginning with code tokens, already delivers appreciable improvements, and further accelerates learning during subsequent math-specific training, cumulatively yielding the strongest results. In contrast, training on a code and math token mixture in a single stage impairs mathematical reasoning when tools are not involved. A plausible explanation is that the parameter count of DeepSeek-LLM 1.3B is insufficient for the model to thoroughly internalize and integrate both code-related and math-related information simultaneously.

\subsubsection{ArXiv Papers Appear Inefficient for Enhancing Mathematical Reasoning}

While ArXiv papers are widely used as a source in math pre-training datasets \citep{minerva,gpt-f,llemma,mathpile}, their specific influence on mathematical reasoning capabilities has not been thoroughly scrutinized. Somewhat unexpectedly, our findings indicate that the inclusion of arXiv content does little to bolster mathematical reasoning performance. We assess models across scales, specifically DeepSeek-LLM 1.3B and DeepSeek-Coder-Base-v1.5 7B \citep{deepseek-coder}, trained using differently preprocessed arXiv datasets:

\begin{itemize}[topsep=0pt]
    \item \textbf{MathPile} \citep{mathpile}: a dataset of 8.9 billion tokens compiled with rigorous cleaning and heuristic filtering, more than 85\% of which are scientific arXiv manuscripts;
    \item \textbf{ArXiv-RedPajama} \citep{redpajama}: a collection comprising all arXiv LaTeX source files, stripped of preambles, comments, macros, and bibliographies, totaling 28.0 billion tokens.
\end{itemize}

For these investigations, DeepSeek-LLM 1.3B is trained for 150B tokens and DeepSeek-Coder-Base-v1.5 7B is trained for 40B tokens on each arXiv-derived corpus independently. Results indicate that arXiv data alone yields negligible gains—and in some cases, even declines—on a spectrum of mathematical reasoning assessments with varying difficulty levels. Such assessments include quantitative problem-solving datasets like GSM8K and MATH (see Table \ref{tab:arxiv-cot}), multiple-choice STEM challenges such as MMLU-STEM (refer to Table \ref{tab:arxiv-cot}), as well as formal mathematics evaluation using miniF2F (Table \ref{tab:arxiv_atp}).

Nevertheless, this finding is not definitive and should be interpreted with caution. Areas left unexplored include:

\begin{itemize}[topsep=0pt]
    \item The role of arXiv data for niche mathematical skills absent from our evaluation, e.g., translating formal mathematical content to informal narratives;
    \item How arXiv data interacts when merged with other datasets;
    \item The potential for arXiv’s utility to emerge at larger model sizes.
\end{itemize}

Consequently, we recognize the need for more comprehensive future research to address these limitations.

\subsection{Insights on Reinforcement Learning}
\subsubsection{Toward a Unified Training Framework} Here, we introduce a general framework that synthesizes various training approaches, including SFT, RFT, DPO, PPO, and GRPO. We also present empirical analysis on key factors within this framework. Broadly, the gradient of the training objective with respect to model parameters $\theta$ for each method can be formulated as:

\begin{equation}
    \nabla_{\theta}\mathcal{J}_{\textcolor{red}{\mathcal{A}}}(\theta) = \mathbb{E}[\underbrace{(q,o) \sim \textcolor{red}{\mathcal{D}}}_{Data \ Source}]\left( \frac{1}{|o|} \sum_{t=1}^{|o|}  \underbrace{GC_{{\mathcal{A}}}(q, o, t, \textcolor{red}{\pi_{{rf}}})}_{Gradient \ Coefficient}  \nabla_{\theta}\log \pi_{\theta}(o_t | q, o_{<t})\right).
\label{eq:objective}
\end{equation}

This framework consists of three fundamental elements: 1) \textit{Data Source $\mathcal{D}$}, specifying which datasets are leveraged during model training; 2) \textit{Reward Function $\pi_{{rf}}$}, supplying the reward signals that drive model updates; and 3) \textit{Algorithm $\mathcal{A}$}, which integrates both the data and reward signals to produce the gradient coefficient $GC$—this value modulates the extent of positive or negative feedback applied during optimization. Within this general structure, we compare several canonical approaches:

\begin{itemize}
    \item  \textbf{Supervised Fine-tuning (SFT)}: This method adapts a pretrained model using labeled datasets curated by humans.
    \item \textbf{Rejection Sampling Fine-tuning (RFT)}: RFT continues to update an SFT model by training on a subset of its outputs; only those responses, generated in response to SFT prompts and verified for correctness, are retained for learning.
    \item \textbf{Direct Preference Optimization (DPO)}: DPO extends SFT through additional tuning on expanded model outputs, optimizing a pairwise objective that prioritizes preferred completions.
    \item \textbf{Online Rejection Sampling Fine-tuning (Online RFT)}: Unlike standard RFT, this strategy initializes the policy network from the SFT model and incrementally fine-tunes it using up-to-date outputs generated by the evolving policy itself.
    \item \textbf{PPO/GRPO}: These algorithms start from an SFT-initialized policy, but employ reinforcement learning, updating the model via samples produced by the live policy network.
\end{itemize}

The components comprising these methods are outlined in Table \ref{tab:data_policy}. For an in-depth derivation and discussion, consult Appendix \ref{app:analysis-rl}.

\paragraph{Discussion of Data Source} We categorize the origins of the training data as either online or offline sampling. Online sampling indicates that data is collected dynamically from the ongoing exploration of the current training policy model. Offline sampling, on the other hand, signifies that data originates from initial extractions using the SFT model. Both RFT and DPO utilize offline-sourced datasets, whereas Online RFT and GRPO leverage online-sampled data.

As depicted in Figure \ref{fig:rl-analysis}, our results reveal that Online RFT achieves considerably stronger performance compared to RFT across two evaluation benchmarks. During the early training phases, Online RFT and RFT yield comparable outcomes; however, as training progresses, Online RFT demonstrates a clear advantage, underscoring the efficacy of online learning approaches. This observation aligns with expectations: at the outset, the behavior of the actor closely matches that of the SFT model, leading to data that reflects only marginal variations. At more advanced stages, though, samples drawn from the actor are increasingly distinct, making online, up-to-date data collection particularly beneficial.

\paragraph{Discussion of Gradient Coefficient} To update model parameters, these algorithms map the reward signal onto a gradient coefficient. In our empirical framework, we distinguish between two types of reward functions: ‘Rule’ and ‘Model’. The ‘Rule’ category involves assessing response quality according to answer correctness, whereas the ‘Model’ category involves employing a trained reward model to assign scores to responses—where this reward model itself was trained using rule-based judgments as supervision. A critical distinction between GRPO and Online RFT, as formalized in Equations \ref{eq:GC-RFT} and \ref{eq:GC-GRPO}, is that GRPO adapts the gradient coefficient dynamically according to the reward model’s output. This property enables the method to reinforce or penalize responses to differing extents, proportional to their scores. Conversely, Online RFT does not possess this adaptive mechanism: it treats all correct responses identically in reinforcement, providing no penalty or differentiation for incorrect answers.

As illustrated in Figure \ref{fig:rl-analysis}, GRPO outperforms online RFT, underlining the effectiveness of modulating positive and negative gradient coefficients. Moreover, GRPO+PS achieves better results than GRPO+OS, which suggests that utilizing finer-grained, step-sensitive gradient adjustments is advantageous. We further investigate the impact of iterative RL by running two iteration cycles in our experiments. The outcomes in Figure \ref{fig:iter-rl} reveal that iterative RL substantially enhances performance, particularly during the initial iteration.

\subsubsection{Why RL Works?} This study applies reinforcement learning on a subset of the instruction tuning dataset, leading to substantial improvements over the base instruction-tuned model. To deepen our understanding of the mechanisms driving these gains, we analyze the Pass@K and Maj@K accuracies for both the Instruct and RL-enhanced models across two benchmarks. The data in Figure \ref{fig:maj-pass} demonstrate that RL mainly boosts the Maj@K metric, while Pass@K remains largely unchanged. This suggests that reinforcement learning fortifies the model's overall output robustness—\textbf{that is, improvements stem from elevating the probability of correct answers among the TopK outputs, rather than upgrading fundamental modeling capabilities}. A related observation by \citep{wang2023making} points to a \textbf{misalignment problem} in reasoning tasks with SFT models, finding that targeted preference alignment techniques can ameliorate reasoning outcomes \citep{yuan2023rrhf,song2023preference,wang2023making}.

\subsubsection{How to Achieve More Effective RL?}
\label{sec:effec_rl}
We show that reinforcement learning performs particularly well on mathematical reasoning tasks. Additionally, we present a unified framework for conceptualizing various representative training strategies, wherein each is treated as a form of direct or approximate RL. According to Equation \ref{eq:objective}, three central elements constitute this paradigm: Data Source, Algorithm, and Reward Function. We outline prospective research opportunities pertaining to each component.

\paragraph{Data Source} The choice of data source forms the foundational basis for all training strategies. In RL, this typically refers to unlabeled problems accompanied by candidate outputs generated by the policy model. Our experiments limit the data source to questions from the instruction tuning phase, with outputs derived using simple nucleus sampling. We hypothesize that this choice may contribute to the observed improvement predominantly in the Maj@K metric. Moving forward, we plan to extend our RL experiments to encompass out-of-distribution prompts and employ \textbf{more sophisticated sampling (decoding) approaches}, such as those that utilize tree-search methods \citep{yao2023tree}. Additionally, the integration of \textbf{efficient inference methods} \citep{xia-etal-2023-speculative,leviathan2023fast,kwon2023efficient,xia2024unlocking}, which critically influence the exploration efficiency of policy models, will be another key avenue of exploration.

\paragraph{Algorithms} Algorithms leverage data and reward signals to adjust the gradient coefficient, thereby updating the model parameters. According to Equation \ref{eq:objective}, contemporary approaches generally \textbf{TRUST} the feedback from the reward function to modulate the conditional probability associated with a specific token, either increasing or decreasing it. Nevertheless, the reliability of reward signals cannot always be guaranteed, particularly for tasks characterized by substantial complexity. As an illustration, even highly curated datasets like PRM800K \citep{lightman2023let}, which have been annotated by experienced professionals, are reported to have around 20\% erroneous labels\footnote{\url{https://github.com/openai/prm800k/issues/12\#issuecomment-1728491852}}. Motivated by this challenge, we aim to investigate reinforcement learning algorithms that demonstrate robustness in the presence of noisy or unreliable reward signals. Our hypothesis is that \textbf{WEAK-TO-STRONG} \citep{burns2023weak} alignment strategies have the potential to transform the foundational mechanisms of current learning algorithms.

\paragraph{Reward Function} The reward function serves as the principal provider of the training signal. Within reinforcement learning, this function is typically implemented as a neural reward model. We identify three pivotal research avenues for reward models: 1) \textbf{Enhancing the generalization capability of reward models.} Effective generalization is crucial so that the reward model can manage out-of-distribution prompts and novel decoding scenarios; if not, reinforcement learning risks merely consolidating the existing model output distribution, without driving genuine progress; 2) \textbf{Quantifying and utilizing the reward model's uncertainty.} Accounting for uncertainty may establish an interface between weak reward models and weak-to-strong learning paradigms; 3) \textbf{Developing high-quality process reward models efficiently} to deliver granular training signals that can more accurately supervise complex reasoning steps \citep{lightman2023let,wang2023math}.

\section{Conclusion, Limitation, and Future Work}

We introduce DeepSeekMath, an open-source model that sets a new benchmark on the MATH competition dataset, outperforming all other open-source contenders and closely rivaling proprietary systems. Built upon DeepSeek-Coder-v1.5 7B, DeepSeekMath was subjected to continual training over 500B tokens, with a substantial subset—120B tokens—comprised of mathematical data from Common Crawl. Through comprehensive ablation experiments, we demonstrate that web pages serve as an outstanding resource for quality mathematical content, whereas arXiv does not provide as much benefit as initially anticipated. We propose Group Relative Policy Optimization (GRPO), a novel approach inspired by Proximal Policy Optimization (PPO), which significantly enhances mathematical reasoning ability with reduced memory requirements. Experimental findings indicate that GRPO remains beneficial for DeepSeekMath-Instruct 7B, even after reaching advanced benchmark scores. Furthermore, we offer a unified framework to contextualize several related techniques and discuss various avenues for advancing reinforcement learning efficacy.

Despite DeepSeekMath's notable performance in quantitative reasoning tasks, its proficiency in domains like geometry and formal theorem proving remains less robust compared to proprietary models. For example, during preliminary assessments, DeepSeekMath struggled with questions involving triangles and ellipses, suggesting that choices made during pre-training and fine-tuning might have introduced certain biases. Additionally, given its model size, DeepSeekMath trails behind GPT-4 in few-shot learning scenarios; whereas GPT-4 benefits from few-shot prompting, DeepSeekMath shows little distinction between its zero-shot and few-shot capabilities. Going forward, we aim to enhance our data selection mechanisms to curate even higher-quality pre-training datasets. Moreover, we will investigate future research directions (as elaborated in Section \ref{sec:effec_rl}) for optimizing reinforcement learning approaches in large language models.

\end{document}